% Môn: Vật lý
% Lớp: 11
% Chương: Chương I. Dao động
% Bài: L11C1 Bài 1. Dao động điều hòa · Khai thác các đại lượng đặc trưng từ đồ thị vận tốc – thời gian hoặc gia tốc – thời gian
% ===== Câu 1 =====
% ID: L11C1B1-236-DS
% Mức: TH
\dangbt{Khai thác các đại lượng đặc trưng từ đồ thị vận tốc – thời gian hoặc gia tốc – thời gian}
\begin{ex}
% ID: L11C1B1-236-DS
% Mức: TH
Vật dao động theo $x=5\cos(3t)$ cm.
\choiceTF
{\True $v=-15\sin(3t)$ cm/s.}
{\True $a=-45\cos(3t)$ cm/s².}
{\True $v_{max}=15$ cm/s.}
{Gia tốc bằng 0 ở biên.}
\end{ex}

% ===== Câu 3 =====
% ID: L11C1B1-237-TLN
% Mức: TH
\begin{ex}
% ID: L11C1B1-237-TLN
% Mức: TH
Với $x=8\cos(2t)$ cm, tốc độ cực đại theo cm/s bằng bao nhiêu?
\shortans{16}
\end{ex}

% ===== Câu 5 =====
% ID: L11C1B1-238-DS
% Mức: VD
\begin{ex}
% ID: L11C1B1-238-DS
% Mức: VD
Vật có $A=6$ cm, $\omega=4$ rad/s, tại $x=A/2$ và đi theo chiều dương.
\choiceTF
{\True Gia tốc âm.}
{\True Độ lớn gia tốc bằng $48\,\text{cm}$ /s².}
{\True Vận tốc dương.}
{\True Vật chậm dần.}
\end{ex}

% ===== Câu 7 =====
% ID: L11C1B1-239-TLN
% Mức: VD
\begin{ex}
% ID: L11C1B1-239-TLN
% Mức: VD
Tại $x=4.5$ cm, $\omega=3$ rad/s. Độ lớn gia tốc theo cm/s² bằng bao nhiêu?
\shortans{40.5}
\end{ex}

% ===== Câu 9 =====
% ID: L11C1B1-240-DS
% Mức: VD
\begin{ex}
% ID: L11C1B1-240-DS
% Mức: VD
Vật đi từ biên âm về vị trí cân bằng.
\choiceTF
{\True Li độ âm.}
{\True Vận tốc dương.}
{\True Gia tốc dương.}
{Tốc độ giảm dần.}
\end{ex}

% ===== Câu 10 =====
% ID: L11C1B1-241-DS
% Mức: TH
\begin{ex}
% ID: L11C1B1-241-DS
% Mức: TH
Từ $x=6\cos(5t+\pi/3)$ cm, lập $v(t),a(t)$ và tính tại $t=0$.
\choiceTF

\end{ex}

% ===== Câu 11 =====
% ID: L11C1B1-242-TLN
% Mức: VD
\begin{ex}
% ID: L11C1B1-242-TLN
% Mức: VD
Vật có $v_{max}=20$ cm/s và $A=5$ cm. Tần số góc theo rad/s bằng bao nhiêu?
\shortans{4}
\end{ex}

% ===== Câu 14 =====
% ID: L11C1B1-243-DS
% Mức: VD
\begin{ex}
% ID: L11C1B1-243-DS
% Mức: VD
Vật có $A=7$ cm, $\omega=2$ rad/s, tại $x=A/2$ và đi theo chiều dương. Tính $v,a$ và xét nhanh dần hay chậm dần.
\choiceTF

\end{ex}

% ===== Câu 18 =====
% ID: L11C1B1-244-DS
% Mức: VD
\begin{ex}
% ID: L11C1B1-244-DS
% Mức: VD
Phân tích dấu của $x,v,a$ và sự thay đổi tốc độ khi vật đi từ biên âm về vị trí cân bằng.
\choiceTF

\end{ex}

% ===== Câu 24 =====
% ID: L11C1B1-245-TN
% Mức: NB
\begin{ex}
% ID: L11C1B1-245-TN
% Mức: NB
Đồ thị quan hệ giữa li độ, vận tốc, gia tốc với thời gian là đường
\choice
{thẳng.}
{elip.}
{parabol.}
{\True hình sin.}
\loigiai{
Trong dao động điều hòa đồ thị quan hệ giữa Li độ, vận tốc, gia tốc đối với thời gian là đường hình sin.
}
\end{ex}

% ===== Câu 28 =====
% ID: L11C1B1-246-TN
% Mức: VD
\begin{ex}
% ID: L11C1B1-246-TN
% Mức: VD
Đồ thị vận tốc – thời gian của một vật dao động điều hòa có vận tốc cực đại $12\,\text{cm}$ /s và chu kì $2\,\text{s}$. Biên độ dao động bằng bao nhiêu?
\choice
{6/π cm}
{6π cm}
{\True 12/π cm}
{12π cm}
\loigiai{
ω = 2π/ $T$ = π rad/s và $A$ = vmax/ω = 12/π cm.
}
\end{ex}

% ===== Câu 32 =====
% ID: L11C1B1-247-TN
% Mức: VD
\begin{ex}
% ID: L11C1B1-247-TN
% Mức: VD
Đồ thị vận tốc – thời gian của một vật dao động điều hòa có vận tốc cực đại $12\,\text{cm}$ /s và chu kì $2\,\text{s}$. Biên độ dao động bằng bao nhiêu?
\choice
{6/π cm}
{6π cm}
{\True 12/π cm}
{12π cm}
\loigiai{
ω = 2π/ $T$ = π rad/s và $A$ = vmax/ω = 12/π cm.
}
\end{ex}

% ===== Câu 36 =====
% ID: L11C1B1-248-TN
% Mức: VD
\begin{ex}
% ID: L11C1B1-248-TN
% Mức: VD
Đồ thị vận tốc – thời gian của một vật dao động điều hòa có vận tốc cực đại $12\,\text{cm}$ /s và chu kì $2\,\text{s}$. Biên độ dao động bằng bao nhiêu?
\choice
{6/π cm}
{6π cm}
{\True 12/π cm}
{12π cm}
\loigiai{
ω = 2π/ $T$ = π rad/s và $A$ = vmax/ω = 12/π cm.
}
\end{ex}

% ===== Câu 40 =====
% ID: L11C1B1-249-TN
% Mức: VD
\begin{ex}
% ID: L11C1B1-249-TN
% Mức: VD
Đồ thị gia tốc – thời gian của một vật dao động điều hòa có gia tốc cực đại $8\,\text{m}$ /s² và chu kì π s. Biên độ dao động bằng bao nhiêu?
\choice
{$1\,\text{m}$}
{\True $2\,\text{m}$}
{$4\,\text{m}$}
{$8\,\text{m}$}
\loigiai{
ω = 2π/$T$ = 2 rad/s. Từ amax = ω²$A$ suy ra $A$ = 8/4 = $2\,\text{m}$.
}
\end{ex}

% ===== Câu 44 =====
% ID: L11C1B1-250-TN
% Mức: VD
\begin{ex}
% ID: L11C1B1-250-TN
% Mức: VD
Đồ thị gia tốc – thời gian của một vật dao động điều hòa có gia tốc cực đại $8\,\text{m}$ /s² và chu kì π s. Biên độ dao động bằng bao nhiêu?
\choice
{$1\,\text{m}$}
{\True $2\,\text{m}$}
{$4\,\text{m}$}
{$8\,\text{m}$}
\loigiai{
ω = 2π/$T$ = 2 rad/s. Từ amax = ω²$A$ suy ra $A$ = 8/4 = $2\,\text{m}$.
}
\end{ex}

% ===== Câu 48 =====
% ID: L11C1B1-251-TN
% Mức: VD
\begin{ex}
% ID: L11C1B1-251-TN
% Mức: VD
Đồ thị gia tốc – thời gian của một vật dao động điều hòa có gia tốc cực đại $8\,\text{m}$ /s² và chu kì π s. Biên độ dao động bằng bao nhiêu?
\choice
{$1\,\text{m}$}
{\True $2\,\text{m}$}
{$4\,\text{m}$}
{$8\,\text{m}$}
\loigiai{
ω = 2π/$T$ = 2 rad/s. Từ amax = ω²$A$ suy ra $A$ = 8/4 = $2\,\text{m}$.
}
\end{ex}
